\bookchapter{الفصل 14: موضوع خاص عن مديري كلمات المرور}

\chapterintro{في العصر الرقمي، كلمات المرور هي مفاتيح باب منزلك. لكن المفاتيح كثيرة ومعقّدة جداً، فما العمل؟ مدير كلمات المرور هو “خادم المفاتيح الشامل” الخاص بك.}

\booksection{أولاً: لماذا تحتاج إلى مدير كلمات المرور؟}

للفرد العادي في العصر الحديث عشرات أو حتى مئات الحسابات على الإنترنت. إذا استخدمت كلمة المرور نفسها لكل حساب، فبمجرد اختراق موقع واحد، تصبح جميع الحسابات في خطر. وإذا استخدمت كلمات مرور معقدة مختلفة لكل حساب، فلا يمكن تذكّرها أصلاً.

يحل مدير كلمات المرور هذه المعضلة: تحتاج فقط إلى تذكّر كلمة مرور رئيسية واحدة، وهو يتذكّر لك جميع كلمات المرور الأخرى.

\booksection{ثانياً: مبدأ عمل مدير كلمات المرور}

\begin{enumerate}
  \item \textbf{التخزين المشفّر}: تُخزَّن جميع كلمات المرور مشفّرة بقوة عالية، ولا يمكن فك تشفيرها إلا بواسطتك.
  \item \textbf{كلمة المرور الرئيسية}: تضبط كلمة مرور رئيسية لفتح خزانة كلمات المرور. لا تُخزَّن كلمة المرور الرئيسية نفسها على الخادم.
  \item \textbf{الملء التلقائي}: في واجهة تسجيل الدخول للمواقع أو التطبيقات، يملأ مدير كلمات المرور اسم المستخدم وكلمة المرور تلقائياً.
  \item \textbf{توليد كلمات المرور}: يمكنه توليد كلمات مرور عشوائية ومعقدة، لرفع مستوى الأمان.
\end{enumerate}

\booksection{ثالثاً: مقدمة عن مبادئ التشفير}

فهم مبادئ تشفير مدير كلمات المرور يساعدك في الحكم على ما إذا كان مزود الخدمة جديراً بالثقة. إليك شرح لبعض المفاهيم الأساسية بلغة مبسطة.

\booksubsection{تشفير AES-256}

AES (معيار التشفير المتقدم) خوارزمية تشفير متماثلة معتمدة من المعهد الوطني الأمريكي للمعايير والتقنية (NIST)، تُستخدم على نطاق واسع في البنوك والجيش والحكومات. “256” تعني طول المفتاح 256 بت — يحتاج نظرياً إلى وقت يتجاوز عمر الكون لكسره بالقوة الغاشمة. مديرو كلمات المرور الرئيسيون (1Password وBitwarden وKeePass وغيرها) يستخدمون AES-2

\booksubsection{بنية المعرفة الصفرية (Zero-Knowledge Architecture)}

“المعرفة الصفرية” لا تعني أن مزود الخدمة لا يخزن شيئاً، بل تعني أن\textbf{التشفير وفك التشفير يتمان محلياً على جهازك، والخادم يحفظ النص المشفر فقط}. العملية المحددة:

\begin{enumerate}
  \item تُدخل كلمة المرور الرئيسية في العميل.
  \item يشتق العميل مفتاح التشفير من كلمة المرور الرئيسية (هذه الخطوة تتم على هاتفك/حاسوبك).
  \item النص المشفر فقط هو ما يُرفع إلى الخادم.
  \item عند تسجيل الدخول التالي، يرسل الخادم النص المشفر، ويفك العميل التشفير بكلمة المرور الرئيسية مرة أخرى.
\end{enumerate}

نقطة رئيسية:\textbf{كلمة المرور الرئيسية نفسها لا تغادر جهازك أبداً}، ولا توجد على الخادم أيضاً. لذا حتى لو اخترق مزود الخدمة، ما يحصل عليه القراصنة مجرد نص مشفر لا يمكن فك تشفيره — ما لم يخمنوا كلمة مرورك الرئيسية.

\begin{tcbinfo}{لماذا تعني "المعرفة الصفرية" أن مزود الخدمة لا يستطيع رؤية كلمات مرورك؟}
لأن مفتاح فك التشفير مُشتق من كلمة مرورك الرئيسية، وهذه الكلمة موجودة في عقلك فقط وتُستخدم على جهازك المحلي فقط. على خادم مزود الخدمة لا توجد كلمة المرور الرئيسية ولا مفتاح فك التشفير، فقط النص المشفر. لهذا يصرح مزودون موثوقون مثل Bitwarden و1Password بوضوح “نحن أيضاً لا نستطيع رؤية كلمات مرورك” — ليس لأنهم لا يريدون، بل لأنهم تقنياً لا يستطيعون.

\end{tcbinfo}

\booksubsection{دالة اشتقاق المفتاح: PBKDF2 و Argon2}

لا يمكن استخدام كلمة المرور الرئيسية مباشرة كمفتاح تشفير — بل تحتاج إلى معالجة عبر “دالة اشتقاق المفتاح”. وظيفة هذه الدالة هي\textbf{جعل الحساب أبطأ عمداً}، لجعل تكلفة الكسر بالقوة الغاشمة على القراصنان مرتفعة لدرجة يصعب معها التنفيذ.

\begin{itemize}
  \item \textbf{PBKDF2}: معيار قديم، يُستخدم على نطاق واسع في iCloud Keychain وBitwarden وLastPass وغيرها. يبطئ الكسر عبر التكرار المتكرر (عادةً أكثر من 100,000 مرة) للتجزئة.
  \item \textbf{Argon2}: بطل مسابقة تجزئة كلمات المرور 2015، أحدث من PBKDF2. بالإضافة إلى "استهلاك القوة الحسابية" يضيف "استهلاك الذاكرة" — مما يجعل تسريع الكسر بالعتاد المتخصص مثل GPU/ASIC صعباً. Bitwarden وKeePassXC وغيرها تدعمه بالفعل.
\end{itemize}

كلما زاد عدد التكرارات زاد الأمان، لكن سرعة فك القفل تتباطأ أيضاً. المزودون الموثوقون يعلنون عن خوارزمية الاشتقاق وعدد التكرارات المستخدمة، لتسهيل تدقيق مجتمع الأمان.

\booksubsection{تحذير من حادثة: تسريب بيانات LastPass 2022}

في أغسطس 2022، أعلن مدير كلمات المرور العريق LastPass عن اختراق بيئة التطوير الخاصة به؛ وفي نوفمبر من العام نفسه كشف عن أن القراصنة استولوا على نسخ احتياطية من خزنة كلمات مرور المستخدمين (نص مشفر). الدروس المستفادة من الحادثة:

\begin{itemize}
  \item \textbf{بنية المعرفة الصفرية حافظت على الحد الأدنى}: بما أن خزنة كلمات المرور مشفرة، حصل القراصنة على النص المشفر لكنهم لم يستطيعوا فك تشفيره — بشرط أن تكون كلمة المرور الرئيسية للمستخدم قوية بما يكفي.
  \item \textbf{أصحاب كلمات المرور الرئيسية الضعيفة هم الأكثر عرضة للخطر}: يمكن للقراصنة كسر كلمات المرور الرئيسية الضعيفة دون اتصال، لذا فإن المستخدمين ذوي كلمات المرور الرئيسية الضعيفة أو المكررة يواجهون خطراً حقيقياً.
  \item \textbf{البيانات الوصفية تسربت أيضاً}: رغم أن محتوى كلمات المرور مشفر، إلا أن البيانات الوصفية مثل عناوين URL للمواقع والملاحظات قد لا تكون مشفرة، مما يكشف الخدمات التي يستخدمها المستخدم.
\end{itemize}

\begin{tcbwarn}{ما تعلمناه من حادثة LastPass}
أولاً، كلمة المرور الرئيسية يجب أن تكون طويلة وفريدة (يُنصح بـ 16 خانة فأكثر، تتضمن أحرف كبيرة وصغيرة وأرقاماً ورموزاً)؛ ثانياً، أعطِ الأولوية للمزودين الذين يدعمون Argon2id وذوي التكرارات العالية؛ ثالثاً، للحسابات عالية القيمة (البنوك، البريد الإلكتروني) فعّل التحقق الثنائي المستقل، ولا تعتمد فقط على مدير كلمات المرور.

\end{tcbwarn}

\booksection{رابعاً: توصيات مديري كلمات المرور الرئيسيين}

\booksubsection{أدوات النظام المدمجة:}

\begin{itemize}
  \item \textbf{iCloud Keychain}: الخيار الأول لمستخدمي منظومة Apple. متكامل بلا انقطاع مع Safari والتطبيقات، مجاني، آمن. العيب هو صعوبة الاستخدام عبر المنصات.
  \item \textbf{Google Password Manager}: مريح لمستخدمي Android. متكامل مع Chrome وAndroid، مجاني.
  \item \textbf{Huawei Password Vault}: مدمج في نظام HarmonyOS، مناسب لمستخدمي Huawei.
\end{itemize}

\booksubsection{أدوات احترافية خارجية:}

\begin{itemize}
  \item \textbf{Bitwarden}: مفتوح المصدر ومجاني، وظائف قوية، يدعم جميع المنصات. أوصي به بشدة للمستخدمين المهتمين بالخصوصية.
  \item \textbf{1Password}: واجهة جميلة، وظائف غنية، مناسب للعائلات والفرق. مدفوع.
  \item \textbf{LastPass}: مدير كلمات مرور قديم، وظائف النسخة المجانية محدودة.
  \item \textbf{KeePass}: تخزين محلي بالكامل، مفتوح المصدر ومجاني. مناسب للمستخدمين المتقدمين.
\end{itemize}

\booksection{خامساً: التقييم الحقيقي لقوة كلمة المرور}

يعرض مديرو كلمات المرور عادةً مؤشر “قوة كلمة المرور” (أحمر/أصفر/أخضر)، لكن هذا التقييم غالباً ما ينظر فقط إلى الطول وأنواع الأحرف، وهو غير شامل. للحكم حقاً على ما إذا كانت كلمة المرور آمنة، يجب مراعاة “الإنتروبيا” و“ما إذا تم تسريبها”.

\booksubsection{الإنتروبيا (Entropy): القياس العلمي لقوة كلمة المرور}

“الإنتروبيا” تقيس عدم يقين كلمة المرور، وحدتها البت. كلما ارتفعت الإنتروبيا، زاد عدد المحاولات اللازمة للكسر بالقوة الغاشمة.

\begin{itemize}
  \item \textbf{"123456"}: إنتروبيا منخفضة جداً، تُكسر في لحظة.
  \item \textbf{"P@ssw0rd"}: تبدو معقدة، لكنها موجودة في كل قوائم هجمات القاموس، والإنتروبيا لا تزال منخفضة.
  \item \textbf{"correct horse battery staple"}: أربع كلمات عشوائية، إنتروبيا حوالي 44 بت، صعوبة الكسر تتجاوز بكثير "P@ssw0rd".
  \item \textbf{16 خانة عشوائية (تتضمن أحرف كبيرة وصغيرة + أرقام + رموز)}: إنتروبيا حوالي 100 بت، لا يمكن كسرها بالقوة الغاشمة بالقوة الحاسوبية الحالية.
\end{itemize}

الخلاصة:\textbf{“الطول + العشوائية” أهم من “الرموز المعقدة”.}. وهذا هو السبب في أن كلمات المرور العشوائية الطويلة التي يولدها مدير كلمات المرور آمنة حقاً.

\booksubsection{HaveIBeenPwned: تحقق مما إذا كانت كلمة مرورك قد تسربت}

الذي يديره الباحث الأمني Troy Hunt\href{https://haveibeenpwned.com}{haveibeenpwned.com}يجمع الحسابات وكلمات المرور التي تسربت في خروقات البيانات الكبرى. يمكنك:

\begin{enumerate}
  \item إدخال بريدك الإلكتروني للتحقق مما إذا كان حسابك قد كُشف في تسريب ما.
  \item استخدام ميزة Pwned Passwords للاستعلام عما إذا كانت كلمة مرور معينة في قاعدة التسريبات المعروفة (مديرو كلمات المرور والمتصفحات دمجوا هذا الاستعلام بالفعل).
\end{enumerate}

إذا ظهرت كلمة مرور في قاعدة التسريبات،\textbf{بغض النظر عن مدى تعقيدها يجب استبدالها فوراً}— لأن هجمات القاموس تجرب هذه الكلمات المعروفة أولاً.

\begin{tcbinfo}{مبدأ استعلام k-إخفاء الهوية}
قلق من أن إدخال كلمة المرور في HaveIBeenPwned قد يتسبب في تسريب آخر؟ لا داعي للقلق. الاستعلامات الحديثة تستخدم بروتوكول “k-إخفاء الهوية”: يرسل المتصفح أول 5 خانات فقط من تجزئة كلمة المرور إلى الخادم، ويعيد الخادم قائمة بكل التجزئات التي تبدأ بهذه الخمس خانات، ثم تتم المقارنة محلياً. أي أن الخادم لا يرى أبداً تجزئة كلمة مرورك الكاملة. 1Password وBitwarden

\end{tcbinfo}

\booksubsection{وظائف المراقبة المدمجة في مدير كلمات المرور}

\begin{itemize}
  \item \textbf{1Password Watchtower}: مراقبة فورية لكلمات المرور الضعيفة والمكررة والمسربة والمنتهية والمواقع التي لا تستخدم 2FA، مع عرض مُصنّف بالألوان أحمر/برتقالي/أصفر.
  \item \textbf{Bitwarden Vault Health}: النسخة المجانية تعرض كلمات المرور الضعيفة والمكررة؛ النسخة المدفوعة تكشف كلمات المرور المسربة (بناءً على HaveIBeenPwned).
  \item \textbf{ميزة "اقتراحات الأمان" في iCloud Keychain}: iOS الإعدادات كلمات المرور اقتراحات الأمان، ينبّهك إلى كلمات المرور التي ظهرت في تسريبات معروفة وتلك المكررة.
  \item \textbf{أداة "فحص كلمات المرور" في Google Password Manager}: مدمج في Chrome، يفحص بكبسة واحدة كلمات المرور المسربة والضعيفة والمكررة.
\end{itemize}

\begin{tcbpractice}{تحقق الآن: هل كلمات مرورك آمنة؟}
\begin{enumerate}
  \item افتح haveibeenpwned.com، أدخل بريدك الإلكتروني المعتاد، وانظر في أي خروقات ظهر.
  \item مستخدمو iPhone: الإعدادات كلمات المرور اقتراحات الأمان، عالج التحذيرات الحمراء واحداً تلو الآخر.
  \item مستخدمو Chrome: الإعدادات كلمات المرور فحص كلمات المرور، عدّل حسب التوجيهات.
  \item مستخدمو 1Password/Bitwarden: افتح Watchtower/Health Reports، أعطِ الأولوية لمعالجة العناصر "المسربة" و"المكررة".
\end{enumerate}

\end{tcbpractice}

\booksection{سادساً: Passkey: مستقبل تسجيل الدخول بدون كلمة مرور}

في 2024-2026، تمر تقنية كلمات المرور بثورة —\textbf{Passkey (مفتاح المرور)}يُتوقع أن يحل محل كلمات المرور التقليدية تماماً. Apple وGoogle وMicrosoft ثلاث منصات كبرى تدفعه بشكل مشترك، بضمان تحالف FIDO واتحاد الويب العالمي (W3C)، وهي أكبر تغيير في طريقة تسجيل الدخول منذ عقود.

\booksubsection{ما هو Passkey؟}

يعتمد Passkey على معايير FIDO2 / WebAuthn، باستخدام\textbf{زوج مفاتيح على الجهاز}بدلاً من كلمة المرور. يمكن تبسيط المبدأ كالتالي:

\begin{itemize}
  \item هاتفك يولّد زوجاً من المفاتيح:\textbf{المفتاح العام}(يُسجل لدى الموقع) +\textbf{المفتاح الخاص}(يبقى دائماً في شريحة الأمان بجهازك).
  \item عند تسجيل الدخول، يرسل الموقع "تحدياً"، يوقّع هاتفك بالمفتاح الخاص، ثم يتحقق الموقع بالمفتاح العام.
  \item قبل التوقيع تحتاج إلى تأكيد هويتك عبر التعرف البيومتري (Face ID / بصمة) أو PIN الجهاز.
  \item في هذه العملية بأكملها،\textbf{لا تُنقل أي كلمة مرور عبر الشبكة}.
\end{itemize}

\booksubsection{Passkey مقابل كلمة المرور التقليدية}

\begin{center}
\begin{tabularx}{\linewidth}{|X|X|X|}
\hline
\textbf{البعد} & \textbf{كلمة المرور التقليدية} & \textbf{Passkey} \tabularnewline
\hline
مقاومة التصيد & ضعيفة، مواقع التصيد يمكن أن تسرق كلمة المرور & مقاومة طبيعية للتصيد، المفتاح الخاص يوقّع على النطاق الحقيقي فقط \tabularnewline
تأثير تسريب الخادم & تجزئة كلمة المرور قد تُكسر & الخادم يخزن المفتاح العام فقط، التسريب بلا تأثير \tabularnewline
السهولة & تحتاج حفظاً/إدخالاً/إدارة & تسجيل دخول بـ Face ID/بصمة في ثانية \tabularnewline
هجوم إعادة التشغيل & قد يُعترض ويعاد تشغيله بواسطة وسيط & كل توقيع لمرة واحدة، لا يمكن إعادة تشغيله \tabularnewline
\hline
\end{tabularx}
\end{center}

\booksubsection{حالة الدعم في كل منصة}

\begin{itemize}
  \item \textbf{Apple iCloud Keychain Passkey}: iOS 16+ / macOS Ventura+ يدعمه أصلاً، يمكن مزامنة Passkey بين جميع الأجهزة التي تسجل دخولها بنفس Apple ID، ولا حاجة لإنشائه من جديد عند تسجيل دخول جهاز جديد للموقع.
  \item \textbf{Google Password Manager Passkey}: Android + Chrome يدعمه أصلاً، يُخزن Passkey في حساب Google، ويتزامن عبر أجهزة Android؛ منذ 2023 يدعم حساب Google الشخصي تسجيل الدخول بـ Passkey.
  \item \textbf{Microsoft}: Windows 11 / Edge يدعم Passkey بالفعل؛ حساب Microsoft وOutlook وOffice وغيرها من الخدمات تفتح تدريجياً تسجيل الدخول بـ Passkey، ويمكن إدارته في Microsoft Authenticator.
  \item \textbf{1Password / Bitwarden}: يدعم إنشاء وتخزين ومزامنة Passkey عبر المنصات، لتبقى غير مقيد بمنظومة واحدة.
\end{itemize}

\begin{tcbinfo}{المزامنة عبر الأجهزة هي الإنجاز الحاسم}
أكبر مشكلة في Passkey مبكراً كانت “الارتباط بجهاز واحد” — تغيير الهاتف يعني إنشاءه من جديد. منذ 2023 نفّذت Apple وGoogle وMicrosoft مزامنة Passkey عبر السحابة: iCloud Keychain يتزامن عبر أجهزة Apple، وGoogle Password Manager يتزامن عبر أجهزة Android. و1Password وغيرها من الأطراف الثالثة ربطت المزامنة عبر المنظومات. هذا يعني أن Passkey أصبح عملياً بما يكفي لاستبدال كلمات المرور.

\end{tcbinfo}

\booksubsection{كيفية إنشاء واستخدام Passkey على الهاتف}

\textbf{مسار iOS (iPhone):}

\begin{enumerate}
  \item في صفحة أمان الحساب على المواقع الداعمة لـ Passkey (مثل Google وMicrosoft وAmazon وeBay وPayPal) اختر "إضافة مفتاح مرور / Add passkey".
  \item ينبثق نظام التحقق Face ID / Touch ID، وبعد اجتيازه يُنشأ Passkey تلقائياً ويُحفظ في iCloud Keychain.
  \item عند تسجيل الدخول التالي للموقع، بعد إدخال اسم المستخدم اختر "تسجيل الدخول بـ iPhone"، مرر Face ID وانتهى.
  \item مدخل الإدارة: الإعدادات كلمات المرور ابحث عن الموقع، ستجد إدخال Passkey.
\end{enumerate}

\textbf{مسار Android:}

\begin{enumerate}
  \item في صفحة أمان الحساب على موقع يدعم Passkey اختر "إضافة مفتاح مرور".
  \item ينبثق Google Password Manager اختر حساب Google تحقق قفل الشاشة (بصمة/نمط/PIN).
  \item يُحفظ Passkey تلقائياً في Google Password Manager ويتزامن عبر أجهزة Android بنفس الحساب.
  \item عند تسجيل الدخول التالي، اختر "استخدام Passkey"، افتح قفل الشاشة وانتهى.
  \item مدخل الإدارة: إعدادات Android كلمات المرور والمفاتيح والتعبئة التلقائية مدير كلمات المرور google.com.
\end{enumerate}

\booksubsection{تسجيل الدخول بـ Passkey عبر الأجهزة}

إذا كنت تسجل الدخول على حاسوب Windows لموقع لديه Passkey على iPhone فقط، يمكنك مسح رمز QR للتفويض:

\begin{enumerate}
  \item في صفحة تسجيل الدخول اختر "استخدام Passkey".
  \item اختر "استخدام جهاز آخر".
  \item يعرض الحاسوب رمز QR، امسحه بـ iPhone.
  \item ينبثق Face ID على iPhone، وبعد التأكيد يكتمل تسجيل الدخول على الحاسوب.
\end{enumerate}

تتم هذه العملية عبر اتصال مشفّر مزدوج بالبلوتوث+الإنترنت، ومفتاح Passkey الخاص لا يغادر iPhone أبداً.

\begin{tcbpractice}{جرّب الآن: جرّب تسجيل الدخول بدون كلمة مرور}
\begin{enumerate}
  \item سجّل الدخول لحساب Google الأمان مفاتيح المرور أضف مفتاح مرور، جرّب تسجيل الدخول بـ Face ID/البصمة.
  \item سجّل الدخول لحساب Microsoft الأمان خيارات الأمان المتقدمة أضف مفتاح مرور.
  \item جرّبه على Taobao وJD وeBay وPayPal وغيرها من الخدمات الداعمة لـ Passkey.
  \item احتفظ بكلمة المرور التقليدية كاحتياط، وانتقل تدريجياً بالحسابات عالية القيمة إلى Passkey.
\end{enumerate}

\end{tcbpractice}

\begin{tcbwarn}{نصيحة لمرحلة الانتقال}
Passkey لم يستبدل كلمات المرور بالكامل بعد — العديد من المواقع ما زالت في مرحلة “Passkey + كلمة مرور” متعايشة. نصيحة: فعّل Passkey أولاً للحسابات عالية القيمة، لكن\textbf{لا تحذف كلمة المرور التقليدية فوراً}، اتركها كاحتياط؛ وتابع استخدام مدير كلمات المرور، لأن ليس كل المواقع تدعم Passkey. اعتبر Passkey “إضافة قوية لمدير كلمات المرور” وليس “استبدالاً فورياً”.

\end{tcbwarn}

\booksection{سبعة: مفاتيح الأمان المادية}

إذا كانت متطلباتك الأمنية أعلى من المستخدم العادي — مثلاً تمتلك عملات مشفرة كبيرة، أو كنت مديراً/صحفياً/ناشطاً، أو تحمي حسابات مؤسسة حساسة — فإن مفتاح الأمان المادي خيار أكثر موثوقية من رمز التحقق عبر الرسائل القصيرة.

\booksubsection{ما هو مفتاح الأمان المادي؟}

المفتاح المادي جهاز مادي (يشبه فلاش USB)، يخزّن شريحة الأمان داخله مفتاحك الخاص، ويحتاج لمساً مادياً أو NFC لإكمال التحقق. منتجات تمثيلية:

\begin{itemize}
  \item \textbf{YubiKey}(شركة Yubico): الأكثر شيوعاً، يدعم FIDO2/WebAuthn وU2F وOTP وPIV وOpenPGP، بواجهات USB-C / USB-A / NFC وغيرها. سلسلة 5 بسعر 50-70 دولاراً.
  \item \textbf{Google Titan}: من إنتاج Google نفسه، USB-C + NFC / بلوتوث، حوالي 30-40 دولاراً.
  \item \textbf{OnlyKey}: بديل مفتوح المصدر، يدعم بروتوكولات متعددة.
  \item \textbf{البديل المحلي الصيني}: مثل Feitian ePasskey، بسعر أكثر اقتصادية.
\end{itemize}

\booksubsection{لماذا المفتاح المادي أكثر أماناً من رمز الرسائل القصيرة؟}

\begin{itemize}
  \item \textbf{مقاومة التصيد}: مثل Passkey، يستجيب المفتاح المادي للنطاق الحقيقي فقط، ومواقع التصيد لا تستطيع خداعه.
  \item \textbf{مقاومة اختطاف SIM}: رمز الرسائل القصيرة يمكن اعتراضه بهجوم SIM-swap، والمفتاح المادي لا يحتاج رقم هاتف.
  \item \textbf{مقاومة رجل في المنتصف}: بروتوكول FIDO2 يتضمن قيمة تحدي في كل توقيع، لا يمكن إعادة تشغيله.
  \item \textbf{عزل مادي}: المفتاح الخاص يبقى دائماً داخل شريحة الأمان، حتى لو اخترق البرمجيات الخبيثة حاسوبك فلن تقرأه.
\end{itemize}

\booksubsection{سيناريوهات الاستخدام}

\begin{itemize}
  \item حسابات منصات تداول العملات المشفرة، واستعادة المحافظ.
  \item SSO المؤسسي، VPN للعمل عن بُعد، تسجيل دخول SSH للخوادم.
  \item البريد الإلكتروني وGitHub ووسائل التواصل الاجتماعي وغيرها من "الحسابات الرئيسية" — بمجرد اختراقها يمكن reset جميع الحسابات الأخرى.
  \item الصحفيون والمحامون والمديرون وغيرهم من أصحاب الحساسية العالية.
\end{itemize}

\booksubsection{كيفية إعداد المفتاح المادي كتحقق ثانوي}

مثال حساب Google + YubiKey:

\begin{enumerate}
  \item اشترِ YubiKey (يُنصح بشراء اثنين دفعة واحدة، واحد للاستخدام اليومي وواحد احتياطي، خزّنهما في أماكن مختلفة).
  \item سجّل الدخول إلى myaccount.google.com الأمان التحقق بخطوتين.
  \item اختر "إضافة مفتاح أمان" أدخل YubiKey في منفذ USB-C أو قرّبه من منطقة استشعار NFC في الهاتف.
  \item اتبع التعليمات والمس النقطة المعدنية في YubiKey لإكمال التسجيل.
  \item يُنصح بتسجيل مفتاح احتياطي ثانٍ.
  \item في إعدادات "التحقق بخطوتين"، اجعل المفتاح المادي طريقة التحقق المفضلة، ورمز الرسائل القصيرة كاحتياطي.
\end{enumerate}

المواقع الرئيسية الداعمة: Google وGitHub وMicrosoft وApple ID وFacebook وTwitter وDropbox وGitLab وAWS وCoinbase وBinance وغيرها.

\begin{tcbinfo}{استراتيجية المفتاح المزدوج}
أكبر خطر للمفتاح المادي هو “الفقدان”. يُنصح: (1) سجّل مفتاحين مختلفين في آن واحد؛ (2) احمل واحداً معك (حلقة المفاتيح)، وخزّن واحداً في خزنة المنزل أو صندوق إيداع البنك؛ (3) بعض المواقع تسمح بطباعة “رموز الاستعادة الاحتياطية” وحفظها دون اتصال. هكذا حتى لو فُقد المفتاح الرئيسي، لن تُحجز خارج حسابك بشكل دائم.

\end{tcbinfo}

\begin{tcbwarn}{احذر من قناة الشراء}
اشترِ حتماً من المتجر الرسمي أو الوكيل المعتمد، ولا تشترِ مفاتيح مستعملة أو من مصدر مجهول — من الناحية النظرية قد تُزرع أبواب خلفية في الأجهزة المعدة مسبقاً. لا يمكن تحديث firmware لـ YubiKey بعد الخروج من المصنع، لكن Yubico يستمر في نشر إعلانات الأمان، تابع موقعه الرسمي.

\end{tcbwarn}

\booksection{ثمانية: مشاركة كلمات المرور للأسرة/الفريق}

في السيناريو العائلي، Netflix وWi-Fi المنزلي ولوحة تحكم الراوتر وكلمة مرور القفل الذكي يحتاجها كل أفراد الأسرة؛ وفي سيناريو الفريق، الخوادم وقواعد البيانات وحسابات التواصل الاجتماعي تحتاج تعاون عدة أشخاص.\textbf{إرسال كلمة المرور عبر WeChat أو كتابتها في مستند مشترك ممارسة عالية الخطورة}— بمجرد تسريب سجل الدردشة أو أذونات المستند، تُسرّب كلمة المرور.

\booksubsection{“الخزنة المشتركة” (Vault) في 1Password}

1Password هو المنتج القياسي لمشاركة الأسرة والفريق. آليته الأساسية هي\textbf{خزنات متعددة}:

\begin{itemize}
  \item \textbf{الخزنة الشخصية}: أنت فقط تراها، مثل الحسابات الخاصة كالبنك والبريد.
  \item \textbf{الخزنة المشتركة}: يراها جميع أفراد الأسرة أو أعضاء الفريق، مثل Netflix وWi-Fi المنزلي والبطاقة الائتمانية المشتركة.
  \item \textbf{تدرج الصلاحيات}: في نسخة الفريق يمكن تعيين أعضاء "عرض فقط" أو "عرض ولكن دون تحرير".
\end{itemize}

لا يحتاج أفراد الأسرة أو أعضاء الفريق لمشاركة كلمة المرور الرئيسية، بل يسجل كل واحد بحسابه الخاص، ويرى محتوى الخزنة المشتركة المصرّح له. تتزامن كلمات المرور فوراً بعد التعديل لجميع الأعضاء.

\booksubsection{وظيفة “المنظمة” في Bitwarden}

النسخة المجانية من Bitwarden تدعم مشاركة كلمات المرور الفردية (عبر وظيفة “إرسال Send” لمشاركة كلمة مرور واحدة محدودة بوقت)؛ النسخة المدفوعة يمكنها إنشاء “منظمة” لمشاركة عدة إدخالات مع أفراد الأسرة أو زملاء الفريق، مع صلاحيات قابلة للتفصيل. يوفر Bitwarden أيضاً خيار الاستضافة الذاتية، مناسب للمؤسسات ذات المتطلبات العالية لسيادة البيانات.

\booksubsection{السيناريو العائلي: طرق آمنة لمشاركة Netflix وWi-Fi وغيرها}

\begin{itemize}
  \item \textbf{Netflix / مواقع الفيديو}: استخدم "الخزنة المشتركة" في 1Password Families لتخزين كلمة مرور تسجيل الدخول، يسجل كل فرد بنفسه دون الحاجة لإرسالها في المجموعة.
  \item \textbf{كلمة مرور Wi-Fi المنزلي}: بين أجهزة iOS يمكن مشاركة كلمة مرور Wi-Fi مباشرة عبر AirDrop للضيوف الزائرين دون الحاجة لنطقها؛ وAndroid 12+ لديه أيضاً وظيفة "مشاركة Wi-Fi القريب".
  \item \textbf{الأقفال الذكية والأجهزة المنزلية الذكية}: استخدم وظيفة "أفراد الأسرة" في تطبيق الشركة لإنشاء حساب منفصل لكل شخص، بدلاً من مشاركة حساب مسؤول واحد.
  \item \textbf{لوحة تحكم الراوتر المنزلي}: خزّن كلمة مرور المسؤول في الخزنة المشتركة لـ 1Password، وافتح SSID منفصل لـ Wi-Fi الزوار.
\end{itemize}

\begin{tcbinfo}{قيمة 1Password Families مقابل السعر}
1Password Families حوالي 4.99 دولار/شهر، يدعم حتى 5 أفراد أسرة، لكل فرد خزنة شخصية مستقلة + 4 خزنات مشتركة. يحسب بأقل من 1 دولار/شهر للفرد، وهو أفضل حل لتجربة المشاركة العائلية حالياً. نسخة Bitwarden العائلية 10 دولار/سنة تدعم 6 أفراد، قيمة أعلى، لكن تجربة الواجهة أقل قليلاً.

\end{tcbinfo}

\booksubsection{سيناريو الفريق: مشاركة كلمات مرور الخوادم/قواعد البيانات}

\begin{itemize}
  \item \textbf{1Password Teams / Business}: يمكن إنشاء خزنات متعددة حسب القسم ومجموعة المشروع، مع تحكم دقيق بالصلاحيات، ويمكن تدقيق جميع سلوكيات الوصول.
  \item \textbf{Bitwarden Teams / Enterprise}: حل مفتوح المصدر، يمكن استضافته ذاتياً، النسخة المؤسسية تدعم تكامل SSO ومزامنة الدليل.
  \item \textbf{مبدأ فصل المسؤوليات}: لا ينبغي أن يكون كلمة مرور المسؤول بيد شخص واحد فقط — عيّن 2-3 "جهات اتصال طوارئ"، لضمان عدم انقطاع العمل عند مغادرة/إجازة الموظفين الرئيسيين.
  \item \textbf{عملية المغادرة}: عند مغادرة الموظف، ألغِ فوراً صلاحية وصوله لمدير كلمات المرور، وأدر كلمات المرور عالية الحساسية التي لمسها عند الضرورة.
\end{itemize}

\booksubsection{الوصول الطارئ (Emergency Access): في حال حدث لك شيء}

تخيل سيناريو: أنت تدير وحيداً جميع الحسابات العائلية المهمة، لكنك دخلت المستشفى فجأة أو توفيّت، ولا يستطيع أهلك تسجيل الدخول للبنك والبريد والتأمين.\textbf{وظيفة “الوصول الطارئ”}صُممت لهذا الغرض — تفويض فرد أسرة/صديق موثوق للحصول على خزنة كلمات المرور عندما لا تستطيع الاستجابة.

\textbf{إعداد الوصول الطارئ في Bitwarden:}

\begin{enumerate}
  \item سجّل الدخول إلى Bitwarden Web Vault الإعدادات الوصول الطارئ.
  \item أضف بريد جهة الاتصال الموثوقة، وادعُه للقبول.
  \item عيّن "فترة الانتظار" (يُنصح بـ 7-14 يوماً): النافذة الزمنية من طلب الطرف الآخر إلى تأكيدك.
  \item إذا لم تستجب خلال فترة الانتظار (ترفض أو لا تفعل شيئاً)، يحصل الطرف الآخر تلقائياً على حق الوصول بعد انتهاء المدة.
  \item يمكنك رفض الطلب يدوياً أو إلغاء صلاحية الوصول الطارئ في أي وقت.
\end{enumerate}

\textbf{الوصول الطارئ في 1Password:}

\begin{enumerate}
  \item 1Password لا يسميه "الوصول الطارئ" مباشرةً، لكن يمكن تحقيقه عبر وظيفة "منظمة الأسرة/الفريق" — بعد إضافة فرد الأسرة للنسخة العائلية، ضع بعض الإدخالات في "الخزنة المشتركة".
  \item لمستخدمي النسخة الشخصية، يُنصح بإخبار فرد أسرة موثوق جداً بكلمة المرور الرئيسية، أو كتابتها في ظرف مختوم يُخزن في الخزنة.
  \item يوفر 1Password أيضاً "رمز الاستعادة"، احرص على طباعته وحفظه عند تفعيل التحقق الثنائي.
\end{enumerate}

\begin{tcbpractice}{اعمله الآن: اترك "مفتاحاً احتياطياً" للأسرة}
\begin{enumerate}
  \item اختر فرد أسرة موثوقاً كـ "جهة اتصال طوارئ".
  \item اضبط الوصول الطارئ في مدير كلمات المرور (Bitwarden) أو أضفه إلى الخزنة المشتركة العائلية (1Password).
  \item اكتب كلمة المرور الرئيسية + رمز الاستعادة على ورقة، ضعها في ظرف، خزّنه في خزنة المنزل أو أخبر فرد أسرة موثوق.
  \item راجع كل نصف سنة، وتأكد ما إذا كانت جهة الاتصال الطارئة لا تزال مناسبة.
\end{enumerate}

\end{tcbpractice}

\booksection{تسعة: أفضل ممارسات استخدام مدير كلمات المرور}

\begin{itemize}
  \item \textbf{اضبط كلمة مرور رئيسية قوية}: كلمة المرور الرئيسية هي الخط الدفاعي الأخير، تأكد من ضبط كلمة مرور طويلة ومعقدة وتذكّرها جيداً.
  \item \textbf{فعّل المصادقة الثنائية}: فعّل المصادقة الثنائية لحساب مدير كلمات المرور، لإضافة طبقة حماية.
  \item \textbf{افحص كلمات المرور الضعيفة دورياً}: لدى معظم مديري كلمات المرور وظيفة "التدقيق الأمني"، تساعدك في العثور على كلمات المرور الضعيفة والمكررة.
  \item \textbf{لا تنسَ كلمة المرور الرئيسية}: لا يمكن استعادة كلمة المرور الرئيسية، ونسيانها يعني فقدان جميع كلمات المرور. يُنصح بكتابتها على ورقة وحفظها في مكان آمن.
  \item \textbf{الوصول الطارئ}: اضبط جهة اتصال للطوارئ، فإذا تعذّر وصولك، يمكن للشخص الموثوق الحصول على كلمات المرور الخاصة بك.
\end{itemize}

\begin{tcbwarn}{تذكير مهم}
مدير كلمات المرور رغم راحته، يحمل مخاطر: إذا تسرّبت كلمة المرور الرئيسية، تنكشف جميع كلمات المرور. لذا يجب ألا تتطابق كلمة المرور الرئيسية مع أي كلمة مرور أخرى، وألا تُكتب في مكان يسهل العثور عليه. كما يُنصح باختيار مزوّد بسجل أمني جيد.

\end{tcbwarn}

\begin{tcbpractice}{جرّب الآن: ابدأ باستخدام مدير كلمات المرور}
\begin{enumerate}
  \item اختر مدير كلمات مرور (يُوصى بـ Bitwarden أو أداة النظام المدمجة).
  \item سجّل حساباً، واضبط كلمة مرور رئيسية قوية.
  \item استورد كلمات المرور الموجودة (إن وُجدت) أو ابدأ بإضافة كلمات مرور جديدة.
  \item ثبّت إضافة المتصفح أو تطبيق الهاتف، وجرّب وظيفة الملء التلقائي.
  \item راجع تقرير التدقيق الأمني، واعثر على كلمات المرور التي تحتاج تحسيناً.
\end{enumerate}

\end{tcbpractice}
